% HAND-WRITTEN fragment -- NOT generated by maestro2tex.
% Source: ~/vestarv/docs/afe_rev2_pixel_verify.html rev. 2026-08-23, S9
% (run simruns/pixtop2_cv, 1.007 s transient, 175 205 stored points, 0 errors):
%   programme  triangle staircase -0.5 -> +0.5 -> -0.5 V vs vcm, 4 mV treads
%              (6-7 DAC codes), 2 ms/tread = 2 V/s, 20 ns bus edges, one cycle,
%              501 treads, codes 1229..2867, ncode = 3;
%   electrode  cdl_we 200 nF, r_u = r_ce = 1 kOhm, e0 0.3 V, k0 0.1, Cox 1e-6,
%              Cred 0, D 5e-6, 37 C -- the same parameter set as the
%              ideal-source run of note_electrode_afe.tex (8.666 uA peak);
%   S9.1  Rf_eff = 37 161.9 ohm median over 124 921 points, p10 36 953.7,
%         p90 37 444.1; nominal law 36 000 ohm; resistor bench 37 620 ohm at 25 C;
%   S9.2  peaks: forward -7.6970 uA @ +0.3436 V (flat bench -8.6658 @ +0.3433),
%         reverse +5.2121 uA @ +0.2551 V (flat bench +6.0553 @ +0.2574);
%         raw (unsampled) peaks are +26 % and +42 % on the same waveform;
%   S9.3  |i_readout - I(IPRB_F)| <= 39.4 nA over 331 treads above 100 nA;
%         Vos trap: the extractor's own pre-scan estimate gives +0.3805 mV.
% Netlist-level in one respect: the 12-bit potential bus is driven by per-bit
% PWL sources so the programme is a real code sequence, and the SAR clock is
% held static.
% 2026-08-24 RE-BASELINE: this 1 s transient was NOT repeated on the edited cell
%   and nothing in simruns/rebase_20260824 moves it.  The two new buffers are
%   outside both loops; deleting the DAC supply regulators left every dc value
%   bit-identical; and stb/stb_table.txt reproduces Zt(1 Hz), peaking and f3dB
%   at every printed digit, so Rf_eff, the peak positions and the readout
%   residual are unaffected.  Re-run it if the loop is ever recompensated.

\subsubsection{Cyclic Voltammetry Through the Complete Chain}\label{sss:pixtop-cv}

Figure~\ref{fig:pixtop-cv} is the electrode of Section~\ref{ss:analog-electrode-bc}
measured through the whole assembled channel: the potential program is a real
12-bit code sequence, the applied potential is servo'd by A1, and the current
is read back through the programmable resistor at gain code~3. The programme is
a triangle staircase, \(\pm\SI{0.5}{\volt}\) about the common mode in
\SI{4}{\milli\volt} treads of 6 or 7 codes, \SI{2}{\milli\second} per tread ---
\SI{2}{\volt\per\second} --- 501 treads over one cycle, sampled once per tread
\SI{50}{\micro\second} before the next code change.

\paragraph{Potential.} The peaks land within \SI{0.3}{\milli\volt} (forward) and
\SI{2.3}{\milli\volt} (reverse) of the same electrode driven by an ideal source.
The potential axis therefore survives \SI{610}{\micro\volt} quantisation, the
ladder's own error, the common-mode DAC's \SI{-262}{\micro\volt} and A1's
\SI{-620}{\micro\volt} offset and still puts the redox features where an ideal
source puts them --- which is what identifies an analyte.

\paragraph{Current.} Sampled peak heights are \SI{11.2}{\percent} and
\SI{13.9}{\percent} below the ideal-source run. A tread sampled at its end has
shed the double-layer charging current that a continuous ramp carries, and the
diffusion layer relaxes over each \SI{2}{\milli\second} hold, so the two
numbers measure different things about the same electrode rather than
disagreeing: kept raw, the same waveform peaks \SI{26}{\percent} and
\SI{42}{\percent} \emph{above} the ideal-source run on the charging tails.

\paragraph{Feedback resistance, measured in the loop.} A current probe in the
feedback branch makes \(R_f\) a measurement rather than an assumption:
\((v_{\mathrm{out}} - v_{\mathrm{we}})/i_f = \SI{37.16}{\kilo\ohm}\) median over
124\,921 points, p10--p90 \SIrange{36.95}{37.44}{\kilo\ohm}, against
\SI{36.0}{\kilo\ohm} from the nominal law and \SI{37.62}{\kilo\ohm} measured
two-terminal at \SI{25}{\celsius} (Section~\ref{ss:tiarprog}) ---
\SI{1.2}{\percent} below the latter at \SI{37}{\celsius}.

\paragraph{Readout.} Over the 331 treads carrying more than
\SI{100}{\nano\ampere} the voltage readout
\((v_{\mathrm{out}} - V_{\mathrm{CM}} - V_{\mathrm{os}})/R_f\) --- the
arithmetic firmware performs --- reproduces the probe current to within
\SI{39.4}{\nano\ampere} on \SI{7.7}{\micro\ampere} peaks, \SI{0.5}{\percent},
inside the \(\pm\SI{0.65}{\percent}\) spread of \(R_f\) itself. No systematic
term is left over. \(V_{\mathrm{os}}\) has to be supplied from a genuine
zero-current measurement (Section~\ref{sss:pixtop-zero}): estimated instead from
the pre-scan quiet period, where this electrode already passes
\SI{27}{\nano\ampere}, it comes out \(+\SI{0.38}{\milli\volt}\) and silently
subtracts a standing current from the whole voltammogram.

One cycle, one gain code, one scan rate, at nominal, with the converter clock
held static.
